\section{Protocol for Detecting Fuzzy Matches}
%
We will describe a protocol the fuzzy matching variant of our content moderation scheme. 
We will assume that the set of blocklist items are embeddings in a metric space $\metricSpace $, and our protocol will detect 
matches of the form $\tupleDefn{\genericBlocklistElmt, \genericCInputElmt}$ such that 
$\distanceNotation[\distance][\genericBlocklistElmt][\genericCInputElmt] \le \threshold$. 
A straightforward idea to extend our base protocol for this scenario is to simply generate a token by 
signing \textit{all} the elements within threshold $\threshold$ of each item in the blocklist. 
However, the number of such items can be 
exponential in $\threshold$. For instance, consider the Hamming space, the number of vectors that 
are within threshold $\threshold$ of some element $\genericBlocklistElmt \in \blocklist$ is 
$\bigO(2^{\threshold})$. For larger values of $\threshold$, we will need exponentially large 
parameters for the signature, e.g., the group order will be super-exponential in 
$\threshold$. 


One possible solution is to use projections of the blocklist items into a low dimensional space for 
comparisons. With careful selection of the projection function, it is possible to 
ensure that matching projections indicate that the corresponding items are close in the metric space 
with some probability better than guessing. On the other hand, since 
projections will introduce false positives, it is possible that two (or more) blocklist elements will 
have the same projection even if they are dissimilar. Now, either we need to find a different 
projection function which produces distinct projections for dissimilar blocklist elements, or 
additional  information needs to be embedded in the blocklist token that will allow us to distinguish 
two colliding blocklist items. For the first approach, we may need several projection functions (linear 
in the blocklist size in the worst case) to distinguish the data, 
while  in the second approach, the information needed can
degenerate to producing a token linear in the size of the blocklist depending on the data 
distribution. 



%We will also allow non-zero false positive and false negatives, i.e.,  if the protocol returns 
%%$\genericSetVar$, then $\genericCInputElmt \in \genericSetVar$ implies
%$\distanceNotation[\distance][\genericCInputElmt.\keyLabel][\genericBlocklistElmt] \le \threshold$ 
%for some $\genericBlocklistElmt \in \blocklist$ with probability at least $\tpr$, \textit{and} 
%$\genericCInputElmt' \notin \genericSetVar$ implies 
%$\distanceNotation[\distance][\genericCInputElmt'.\keyLabel][\genericBlocklistElmt] > 
%\threshold$ for all $\genericBlocklistElmt \in \blocklist$ with probability at least $\tnr$. 

To solve these problems, our protocol will utilize a general primitive called a 
\projectionPrimitiveIt, which is implied by a construction due to 
\citet{chongchitmate2024approximate}. We formalize the primitive below 

\begin{definition}
	A \projectionPrimitive of a set $\blocklist \subseteq \metricSpace$ 
	comprises the projections of each constituent item under uniformly sampled functions $\projFn[1], 
	\dots, \projFn[\numProjections]: \metricSpace \rightarrow \projectedSpace$ 
	
	\begin{enumerate}
		\item Given any two items 
		$\genericMetricSpacElmt,\genericMetricSpacElmtAlt \in \metricSpace$
		
		
		\begin{align}
			\distanceNotation[\distance][\genericMetricSpacElmt][\genericMetricSpacElmtAlt] \le 
			\threshold,  \prob{\forall \projIdx \in \nats[1][\numProjections]: 
				\projFn[\projIdx](\genericMetricSpacElmt) = \projFn[\projIdx](\genericMetricSpacElmtAlt)} 
				\ge 
			1 - \epsilon \label{eq:no_false_neg} \\
			\distanceNotation[\distance][\genericMetricSpacElmt][\genericMetricSpacElmtAlt] > 
			\threshold, \prob{\forall \projIdx \in 
				\nats[1][\numProjections]: \projFn[\projIdx](\genericMetricSpacElmt) \neq
				\projFn[\projIdx](\genericMetricSpacElmtAlt)} \ge 1 - \epsilon \label{eq:no_false_pos}
		\end{align}
		
		%Here, \tpr and \tnr are tunable true positive rate and true negative rate. 
		
		\item For every element 
		$\genericBlocklistElmt \in \blocklist$, there is at least one function $\projFn[\projIdx]$ where 
		$\projFn[\projIdx](\genericBlocklistElmt) \notin 
		\setDefn{\projFn[\projIdx](\genericBlocklistElmt[1]'), 
			\dots, \projFn[\projIdx](\genericBlocklistElmt[\blocklistSize]')}$ for all 
		$\genericBlocklistElmt[1]', 
		\dots, \genericBlocklistElmt[\blocklistSize]' \in \blocklist\setminus\genericBlocklistElmt$ (with 
		overwhelming probability).
		
	\end{enumerate}
\end{definition}

Intuitively, the first condition ensures that items close in the metric space have the same projections 
with high probability. Conversely, the items that are far apart do not agree 
on any of the projections with high probability. Note that $\epsilon$ may not be negligible, and 
rather it is useful to think of it as a non-zero false positive and/or false negative rate. 
The second condition implies that the 
projections partition the items in \blocklist such that there is at least one unique projection for each 
constituent item. In particular, consider the following algorithm that partitions the set $\blocklist$
into $\numProjections$ buckets $\projBuckets[1], 
\dots, \projBuckets[\numProjections]$. The non-colliding projections of items in \blocklist are added 
to the buckets sequentially, i.e., the projections of all items that do not collide under $\projFn[1]$ are 
added to $\projBuckets[1]$. Then, the projections of all items that do not collide under $\projFn[2]$ 
are added to $\projBuckets[2]$. This process 
is repeated for all \numProjections functions. It is guaranteed that every item in $\blocklist$ has 
at least one projection in 
$\bigcup\limits_{\projIdx \in \nats[\numProjections]} \projBuckets[\projIdx]$ with high 
probability\footnote{If we we assume that the blocklist only contains pairwise-dissimilar items, then 
the first condition may imply the second condition for small security 
parameters~\cite{chongchitmate2024approximate}}. 

\myparagraph{Construction sketch}
%
A \projectionPrimitive of the auditors set $\blocklist$ gives a straightforward mechanism to extend 
our base protocol to detect fuzzy matches. The auditors partition $\blocklist$ into sets 
$\projBuckets[1], \dots, \projBuckets[\numProjections]$ as described above. Then, they  
individually sign $\projBuckets[1], \dots, \projBuckets[\numProjections]$ using a randomizable 
set homomorphic signature and publish $\genericSig[\projBuckets[1]], \dots, 
\genericSig[\projBuckets[\numProjections]]$. The client projects its inputs 
$\genericCInputElmt$ (the content embedding)
using $\projFn[1], \dots, \projFn[\numProjections]$ and then computes 
$\genericSig[\projBuckets[1]] \sigMinus[\pubKey]\projFn[1](\genericCInputElmt)$, $\dots$, 
$\genericSig[\projBuckets[\numProjections]] 
\sigMinus[\pubKey]\projFn[\numProjections](\genericCInputElmt)$
(see \stepsref{nf:client:sigminus}{nf:client:randomize} of the non fuzzy protocol). 
The client encrypts the content under the key $\genericSymmKey$ and then encrypts 
\genericSymmKey under $\numProjections$ different keys 
producing $\numProjections$ ciphertexts: the $\projIdx$-th key is derived from the output of 
$\randomize[\pubKey](\genericSig[\projBuckets[\projIdx]] 
\sigMinus[\pubKey]\projFn[\projIdx](\genericCInputElmt))$
(\stepref{nf:client:encrypt}). Upon receiving the ciphertexts and the signatures, the server tries to 
decrypt the ciphertexts (\stepsref{nf:server:check}{nf:server:decrypt}).  The decryption is 
successful if the server is able to correctly compute one of the keys, which happens when there is 
some $\genericBlocklistElmt \in \blocklist$ that has the same projection as $\genericCInputElmt$ 
under at least one out of $\projFn[1], \dots, \projFn[\numProjections]$. 


The correctness of the 
protocol follows from the fact for all $\genericBlocklistElmt \in \blocklist$, there is at least one 
projection $\projFn[\projIdx]$ such that $\projFn[\projIdx](\genericBlocklistElmt) \in 
\projBuckets[\projIdx]$ and if 
$\distanceNotation[\distance][\genericBlocklistElmt][\genericCInputElmt] \le \threshold$, 
then $\projFn[\projIdx](\genericBlocklistElmt) = \projFn[\projIdx](\genericCInputElmt)$ with 
probability at least $1 - \epsilon$. Thus, from 
$\genericSig[\projBuckets[\projIdx]]\sigMinus[\pubKey]\projFn[\projIdx](\genericCInputElmt)$,
the server will be able to derive the $\projIdx$-th key to decrypt $\genericCInputElmt.\valueLabel$. 
The security of  the framework is implied by the signature scheme: if 
$\projFn[\projIdx](\genericCInputElmt) \notin \projBuckets[\projIdx]$ for all $\projIdx \in 
\nats[\numProjections]$, then 
$\genericSig[\projBuckets[\projIdx]]\sigMinus[\pubKey] 
\projFn[\projIdx](\genericCInputElmt)$ provides no information about 
$\projFn[\projIdx](\genericCInputElmt)$ and any of the keys 
encrypting the content to the server due to the IND-SR property. Also, 
$\projFn[\projIdx](\genericCInputElmt) \neq \projFn[\projIdx](\genericBlocklistElmt)$ for all 
$\projIdx \in \nats[\numProjections]$ with probability at least $1 - \epsilon$ when  
$\distanceNotation[\distance][\genericBlocklistElmt][\genericCInputElmt] > \threshold$. 
Finally, observe that the client needs to compute and communicate $\bigO(\numProjections)$ 
elements in the protocol and the server's compute cost is $\bigO(\numProjections \blocklistSize 
\log\blocklistSize)$. 


\myparagraph{Candidate \projectionPrimitive schemes}
%
Before describing a concrete instantiation of the sketch, it 
is worth discussing the candidates we can use for a \projectionPrimitive.
The obvious candidate is a locality sensitive hash function $\localitySensitiveHash: \metricSpace 
\rightarrow \projectedSpace$ which (informally stated) guarantees that for any two elements 
$\genericMetricSpacElmt, \genericMetricSpacElmtAlt \in 
\metricSpace$, if 
$\distanceNotation[\distance][\genericMetricSpacElmt][\genericMetricSpacElmtAlt] \le \threshold$ 
then $\prob{\localitySensitiveHash(\genericMetricSpacElmt) = 
	\localitySensitiveHash(\genericMetricSpacElmtAlt)} \ge \tpr$, where \tpr is the true positive rate. 
Conversely, if 
$\distanceNotation[\distance][\genericMetricSpacElmt][\genericMetricSpacElmtAlt] > 
\thresholdUpper$, then $\prob{\localitySensitiveHash(\genericMetricSpacElmt) \neq
	\localitySensitiveHash(\genericMetricSpacElmtAlt)} \ge \tnr$. If $\thresholdUpper \ge 
\lshConst \threshold$ for some integer constant $\lshConst > 1$, then by concatenating the 
outputs of \numProjections different uniformly sampled hash functions (which are essentially 
projections), false positives and false 
negatives can be eliminated with overwhelming probability~\cite{}.
It is not immediately obvious that locality-sensitive hash functions will ensure condition (2), 
and indeed without making additional assumptions about the data distribution, it
seems unlikely that they will. However, under certain distribution assumptions (about the items in 
the blocklist), we can achieve this guarantee~\cite{chongchitmate2024approximate}. The 
assumption roughly translates to: if there are two dissimilar items in the blocklist, then they are 
sufficiently far apart so that a sublinear (in the input set size) number of random projections are 
enough to distinguish them; this guarantee is formalized in \secref{}. Crucially, we do not need to 
make any assumptions about the client's input, and their relation to the blocklist items\footnote{An 
	alternative solution may also be constructed from spatial hashing techniques introduced 
	by~\citet{} 
	but we defer this discussion to future works}.

While making an assumption about the blocklist might appear limiting, there are several 
reasons why they are justified for content moderation. First, many embedding 
functions implicitly guarantee that dissimilar items will be embedded far apart in the metric space. A 
noteworthy example are perceptual hash functions that map images to Hamming vectors, and are 
used in state of the art CSAM detection systems~\cite{}. Perceptual hash functions ensure that for 
two dissimilar images, the normalized Hamming distance is roughly 0.5, which is expected the 
Hamming distance between any two random images (see \cite{} for a comparison of different 
perceptual hash techniques). Second, the auditors can decide the items to include in the 
blocklist, and the type of embedding function. If an embedding results in a large number of 
collisions, it can be tuned/replaced by a different function. Finally, if such a partitioning is not 
possible for all items, it is possible to simply revert to a mechanism that does not require such a 
partitioning, e.g., a fully homomorphic encryption based 
scheme~\cite{}, for a small subset of the blocklist. 





\subsection{Instantiation of a Delegated Fuzzy PSI}
%
In this section, we will describe a protocol that works with blocklist containing content embeddings 
in the Hamming space. More specifically, assume that there is a content embedding 
function $\universe \rightarrow \setDefn{0,1}^{\genericVecLen}$ and a similarity metric 
$\similarityMetric: \universe \times \universe \rightarrow \setDefn{0,1}$, which given 
$\genericUniverseElmt, \genericUniverseElmtAlt \in \universe$, returns 1 if the content is considered 
similar to each other, and returns 0 otherwise. The embedding function ensures for $\threshold \in 
\nonnegativeIntegers$

\itemListStart
\item 
$\cprob{\big}{\hammingDist{\contentHash(\genericUniverseElmt)}{\contentHash(\genericUniverseElmtAlt)}
	\le 
	\threshold}{\similarityMetric(\genericUniverseElmt, \genericUniverseElmtAlt) = 1} \ge \tpr$
\item 
$\cprob{\big}{\hammingDist{\contentHash(\genericUniverseElmt)}{\contentHash(\genericUniverseElmtAlt)}
	> 
	\threshold}{\similarityMetric(\genericUniverseElmt, \genericUniverseElmtAlt) = 0} \ge \tnr$
\itemListEnd

, where $\hammingDist{\cdot}$ is the Hamming distance between two vectors, \tpr is the true 
positive rate and \tnr is the true negative rate of the embedding function. These parameters can be 
tuned for embedding functions by e.g., adjusting the output vector lengths, etc. Many common 
metric distances, e.g., $L_p$, Edit distance, etc. can be embedded into Hamming 
distance using known techniques~\cite{}. 

The auditors first maps the blocklisted content to bit-vectors using $\embed$ and creates the 
embedded set $\blocklistEmbedded$. This set is then partitioned using a 
\projectionPrimitive scheme over Hamming vectors to create the projected sets $\projBuckets[1], 
\dots, \projBuckets[\numProjections]$. We will leverage a construction by 
\citet{chongchitmate2024approximate} as a building block 

\myparagraph{\ProjectionPrimitive of Hamming 
	vectors~\cite{chongchitmate2024approximate}}
%
The construction uses a set of uniformly sampled functions $\projFn[1], \dots, 
\projFn[\numProjections]$ where $\projFn[\projIdx]: \setDefn{0,1}^{\genericVecLen} \rightarrow 
\setDefn{0,1}^{\secParam}$. Each such function selects a random subset of bits from the input 
vector, and injectively maps it to a $\secParam$-bit string. Consider a set of bit-vectors 
$\blocklistEmbedded \subseteq \setDefn{0,1}^{\genericVecLen}$. For each function 
$\projFn[\projIdx]$ and each element $\genericVec \in \blocklistEmbedded$, we evaluate 
$\projFn[\projIdx](\genericVec)$; the function evaluation is essentially a projection 
of $\genericVec$. Denote the set of non-colliding projections under $\projFn[\projIdx]$ as
$\projBuckets[\projIdx]$, then $\projBuckets[1], \dots, \projBuckets[\numProjections]$ is a 
\projectionPrimitive of \blocklistEmbedded for an appropriate value of $\numProjections$. 
Specifically, if all the vectors in $\blocklistEmbedded$ are at least $\thresholdUpper = \lshConst 
\threshold$ Hamming distance apart, then $\numProjections = \bigO((\genericSetSize 
\lshConst)^{\frac{1}{\lshConst - 1}} \secParam)$ functions partition the items in 
$\blocklistEmbedded$ with overwhelming probability (failing negligibly in 
$\secParam$). Observe that if $\lshConst > 2$, the number of functions required is sublinear in 
$\setSize{\blocklistEmbedded} = \genericSetSize$. Most embedding functions ensure that similar 
items are mapped to items close in the embedded space, while dissimilar items are mapped to items 
that are sufficiently far apart~\cite{}. Therefore if $\blocklist$ contains dissimilar items (which is 
likely since similar items can be clustered together), $\numProjections$ is sublinear in 
$\setSize{\blocklist}$. 



%Consequently, if the embedding 
%function introduces false negatives, i.e., cases where 
%$\distanceNotation[\distance][\genericBlocklistElmt][\genericCInputElmt] \le \threshold$ 
%but $\hammingDist{\embed(\genericBlocklistElmt)}{\embed(\genericCInputElmt.\keyLabel)} >
%\scalingFn(\threshold)$ with false negative rate $1 - \tpr$,  then combining it with the 
%\projectionPrimitive described above ensures an overall construction with identical false 
%negative rate. Since \tpr of the embedding function is tunable, e.g., by varying the 
%length of the vectors, the false negative rate of the overall scheme is also tunable using 
%the same parameters. 

\myparagraph{Towards a protocol design}
%
There is one technical challenge that we need to address before using the partitioning mechanism 
in our construction sketch. First, while the schemes ensures that false negatives 
are negligible (\eqnref{eq:no_false_neg}), it does not guarantee the same for false positives.
That is, even when $\hammingDist{\genericBlocklistElmt}{\genericCInputElmt} >
\threshold$, it is possible that   
$\projFn[\projIdx](\genericBlocklistElmt) 
= \projFn[\projIdx](\genericCInputElmt)$ for some $\projIdx \in \nats[\numProjections]$. These false 
positives will
compound the false positives from the embedding function, thereby reducing the accuracy of 
detection. One solution is to additionally check every matched pair and eliminate false 
positives. When implemented privately in a two-party setting (e.g., 
\cite{chongchitmate2024approximate}), eliminating false positive requires an additional 
protocol that computes the Hamming distance between pairs with matching 
projections. 


However, a similar fine-grained check is not immediately compatible with our signature scheme.  
Specifically, to perform this check, the auditors will need to sign the individual components of each 
vector in the embedded blocklist. Therefore, the signature will be over a multi-set. 
For example, consider the embedded blocklist with two vectors 
$110$ and $100$. The multi-set includes $\{(1,1), (1,1), (1,2), (0,2), (0,3), (0,3)\}$, assuming that 
each bit value is also tagged with its corresponding component. Now, since we accumulate the set 
before signing, additional information must be included in the signature that will allow us to 
distinguish conflicting elements, e.g., a token indicating whether the element $(1,1)$ is part of the 
first or second vector. Moreover, if a client's input vector is say $111$, it would also need to generate 
metadata allowing the server to  check $(1,1)$ against both the token values. As a result the 
metadata size no longer remains sublinear in the blocklist size. 




\myparagraph{Protocol intuition:}
%
We present a solution that subsumes the 
need for an additional protocol. Intuitively, the idea is that we 
will tie the projections of a given vector with its bit-values e.g., using a collision-resistant hash 
function, and accumulate over the outputs. According to this scheme, two vectors will agree on a 
particular component if: i) they have the same bit value at that component, and ii) they have the 
same projection under at least one of the projection functions. Using this criteria, we will compare 
the vectors component-wise. Following the previous example, the set that the auditors will need to 
sign is $\{\hashFn(\projIdx, 1, 1, \projFn[\projIdx](110)), \hashFn(\projIdx, 1, 2, 
\projFn[\projIdx](110)), \\ \hashFn(\projIdx, 0, 3, 
\projFn[\projIdx](110)), \hashFn(\projIdx, 1, 1, \projFn[\projIdx](100)),$ $\hashFn(\projIdx, 0, 2, 
\projFn[\projIdx](100)), 
\hashFn(\projIdx, 0, 3, \projFn[\projIdx](100))\}$. Note that this is \textit{not} a multi-set assuming 
that the hash function is collision resistant and $\projFn[\projIdx](110) \neq \projFn[\projIdx](100)$ 
due to condition (2) of a \projectionPrimitive. Also note that we are including the projection number 
in the hash because we want to avoid collisions across different projection functions. 

More generally, the set of embedded bit-vectors $\blocklistEmbedded$ are first partitioned into the 
buckets $\projBuckets[1], \dots, \projBuckets[\numProjections]$ as described before 
(\stepref{auditor:partition} of the protocol for detecting fuzzy matches). The 
bucket $\projBuckets[\projIdx]$ contains the projection $\projFn[\projIdx](\genericVec)$
for each bit-vector $\genericVec = \tupleDefn{\genericVec[1], \dots, \genericVec[\genericVecLen]} 
\in \blocklistEmbedded$.  Then, we create an augmented set 
$\projBucketsAugmented[\projIdx]$ where corresponding to each entry 
$\projFn[\projIdx](\genericVec)$
in $\projBuckets[\projIdx]$, we include the items $\hashFn(\projIdx, \genericVec[1], 1
\projFn[\projIdx](\genericVec)), \dots, \hashFn(\projIdx, \genericVec[\genericVecLen], 
\genericVecLen, \projFn[\projIdx](\genericVec))$ (\stepref{auditor:augment}) where $\hashFn$ is a 
collision resistant hash function. The auditors sign $\projBucketsAugmented[1], \dots, 
\projBucketsAugmented[\numProjections]$
using a randomizable set-homomorphic signature scheme and publish the signatures 
(\stepref{auditor:sign})\footnote{Our protocol uses the set homomorphic signature construction 
described 
in \secref{sec:sig_scheme} and thus subsumes the need for a separate collision 
resistant hash function because items are mapped to primes using a random oracle before being 
signed; nonetheless we describe the idea more generally}.

During an upload, let the client's input after 
embedding into a bit-vector be $\clientVec = \tupleDefn{\clientVec[1], \dots, 
	\clientVec[\genericVecLen]}$. Then the 
client computes the set of elements 
$\setDefn{\randomize[\pubKey](\genericSig[\projBucketsAugmented[\projIdx]]) 
	\sigMinus[\pubKey]\hashFn(\projIdx, \clientVec[\vecIdx],
	\vecIdx, \projFn[\projIdx](\clientVec[\vecIdx]))}_{\vecIdx \in \nats[\genericVecLen]}$ for $\projIdx 
\in \nats[1][\numProjections]$ (\stepref{client:randomize}). Note that  if 
$\hashFn(\projIdx, \clientVec[\vecIdx],
\vecIdx, \projFn[\projIdx](\clientVec[\vecIdx])) \in \projBucketsAugmented[\projIdx]$, then there is 
some vector $\genericVec 	= \tupleDefn{\genericVec[1], \dots, \genericVec[\genericVecLen]} \in 
\blocklistEmbedded$ such that $\clientVec$ and $\genericVec$ have the same projection under 
$\projFn[\projIdx]$ \textit{and} $\clientVec[\vecIdx] = \genericVec[\vecIdx]$ with overwhelming 
probability. Also, observe that if $\hashFn(\projIdx, \clientVec[\vecIdx],
\vecIdx, \projFn[\projIdx](\clientVec[\vecIdx])) \in \projBucketsAugmented[\projIdx]$ for at least 
$\genericVecLen - \threshold$ different values of $\vecIdx$, then 
$\hammingDist{\clientVec}{\genericVec} \le \threshold$. 
This follows from the fact that $\projFn[\projIdx](\genericVec)$ is not equal to the projection of any 
other vector in $\blocklistEmbedded$ under $\projFn[\projIdx]$, thus the set of values 
$\setDefn{\hashFn(\projIdx, \genericVec[1], 
	1, \projFn[\projIdx](\genericVec)), \dots, \hashFn(\projIdx, \genericVec[\genericVecLen], 
	\genericVecLen, 
	\projFn[\projIdx](\genericVec))}$ is unique to $\genericVec$ assuming that $\hashFn$ is collision 
resistant. 

\myparagraph{Detecting threshold matches}
%
Similar to the non-fuzzy protocol, the client will use 
$\setDefn{\randomize[\pubKey](\genericSig[\projBucketsAugmented[\projIdx]]\sigMinus[\pubKey] 
	\hashFn(\projIdx, \clientVec[\vecIdx],
	\vecIdx, \projFn[\projIdx](\clientVec[\vecIdx]))) = \sigPair[\projIdx, \vecIdx]}$ to encrypt
$\clientData$. More specifically, we want to ensure that the server is able to 
decrypt the content \textit{only} when it can compute $\sigPair[\projIdx, \vecIdx][2]$ for at least 
$\genericVecLen - \threshold$ values of $\vecIdx \in \nats[\genericVecLen]$.
For this, the client secret shares the key using a $(\genericVecLen - 
\threshold)$-out-of-$\genericVecLen$ secret sharing scheme, and encrypts the shares 
using $\sigPair[\projIdx,\vecIdx][2]$ as a one-time pad for $\vecIdx \in \nats[\genericVecLen]$  
(\stepsref{client:share}{client:dhf}); when using the set-homomorphic scheme described 
in \secref{sec:sig_scheme}, this is trivial since $\sigPair[\projIdx,\vecIdx][2]$ is the output of a 
random oracle. Upon receiving the signatures  $\sigPair[\projIdx,\vecIdx][1]$ for 
all $\vecIdx \in \nats[\genericVecLen]$ and the encrypted secret shares, the server tries to compute 
$\sigPair[\projIdx,\vecIdx][2]$ for at least $\genericVecLen - \threshold$ values, 
which gives it access to the required number of shares to compute $\genericSymmKey$ and decrypt 
$\encrypt[\genericSymmKey](\clientData)$ (\stepref{server:verify}). 

The remaining challenge is that the server needs to know which out of the $\genericVecLen$ shares 
have been correctly computed. The caveat here is that if less than  $\genericVecLen - 
\threshold$ shares have been correctly computed -- which implies that with overwhelming 
probability $\forall \genericVec \in \blocklistEmbedded, \hammingDist{\genericVec}{\clientVec} > 
\threshold$ -- then we do not want to reveal the \textit{locations} of the correctly computed shares. 
Note that the shares the server obtains correctly indicates the components where the 
vectors $\clientVec$ and $\genericVec \in \blocklistEmbedded$ have the same value. 
Since the embedding function maps close items 
in $\universe$ to vectors with low pairwise Hamming distance, revealing these components will 
enable the server to learn more information about the client's data, even when it is not able to 
decrypt the content.  


We solve this problem using a detectable hash function. Specifically, along with the secret share, 
we also encrypt (using the one-time pad) the output of the DHF on a random value. If the server 
correctly computes the random pad, then it not only obtains the secret-share of the key, but also a 
DHF output (\stepref{client:dhf} and \stepsref{server:dhf}{server:decrypt}). Consequently, with at 
least $\genericVecLen - \threshold$ DHF outputs, the 
server learns which $\threshold$ out of the $\genericVecLen$ shares have been correctly 
decrypted. Otherwise, the DHF 
outputs and the decrypted shares look random to the server. 



\setcounter{auditorLineNmbr}{0}
\setcounter{clientLineNmbr}{0}
\setcounter{serverLineNmbr}{0}

\begin{mdframed}[frametitle={Protocol for Detecting Fuzzy Blocklist Matches}, 
frametitlealignment=\centering, 
font=\small]

	\smallskip\noindent\textbf{Parameters:} A randomizable set homomorphic signature 
	scheme 
$\sigScheme \assign \genericSigSchemeDefn$ with signature space $\sigOutputSpace$; an 
IND-CPA secure encryption scheme $\encryptionSchemeDefn$ with key space 
$\encryptKeySpace$; set of uniformly sampled functions $\projFn[1], \dots, 
\projFn[\numProjections]: \setDefn{0,1}^{\genericVecLen} \rightarrow \projRange$ for a 
locality-preserving partition over Hamming vectors; an embedding function 
$\contentHash: \universe \rightarrow \setDefn{0,1}^{\genericVecLen}$; 
random oracles $\primeRO: 
\setDefn{0,1} \times \nats[\genericVecLen] \times \projRange \rightarrow \primes$;  
a $(\threshold, \genericVecLen - \threshold)$ detectable hash 
function $\dhf: \dhfKeySpace \times \dhfDomain \rightarrow \dhfRange$;
a secret-sharing scheme $\secretSharingScheme: \encryptKeySpace \rightarrow 
\secretShareOutputSpace^{\genericVecLen}$; random oracle $\ro: 
\sigOutputSpace \rightarrow \setDefn{0,1}^{\setSize{\dhfRange} + 
	\setSize{\secretShareOutputSpace}}$;

\smallskip\noindent\textbf{Inputs:}


\begin{itemize}[leftmargin=0.5em]
\item The auditor \auditor has a blocklist $\blocklistEmbedded$ with \blocklistSize content 
embeddings derived using \contentHash, 
and the public-private key pair for the signature scheme $\pubKey, \privKey$. Let $\pubKey 
\assign \tupleDefn{\rsaMod, \qrGen}$ where $\rsaMod = 
\safePrime[1]  \safePrime[2]$ and $\privKey \assign \safePrime[1], \safePrime[2]$ (see 
\secref{sec:sig_scheme})

\item The client \client has the public key $\pubKey$ and data $\clientData \in 
\setDefn{0,1}^{\ast}$ that it wants to upload to the server. 

\item The server \server has the public key $\pubKey$ and the embedded blocklists 
\end{itemize}	

\smallskip\noindent\textbf{Auditor} ($\createBlocklist(\privKey, \blocklistEmbedded)$)
%


\begin{enumerate}[leftmargin=1.5em]

	\item[\auditorLabel{auditor:selectPrime}] $\blindingFactor_1, \dots, 
	\blindingFactor_{\numProjections} \getsr \primes$ 

\item[\auditorNoLabel] $\forall \projIdx \in \nats[\numProjections], 
\projBuckets[\projIdx] \assign \emptyset$
\item[\auditorLabel{auditor:partition}] $\forall \projIdx \in \nats[\numProjections], \forall 
\genericVec \in \blocklistEmbedded, 
\projBuckets[\projIdx] =  \projBuckets[\projIdx] \cup \projFn[\projIdx](\genericVec) 
~\mbox{if}~ 
\projFn[\projIdx](\genericVec) \notin \projBuckets[\projIdx]$ 
\item[\auditorLabel{auditor:augment}] $\forall \projIdx \in \nats[\numProjections],
\projBucketsAugmented[\projIdx] \assign \setDefn{\primeRO(\projIdx, \genericVec[\vecIdx], \vecIdx, 
	\projFn[\projIdx](\genericVec))}_{\genericVec \in \blocklistEmbedded, \vecIdx \in 
	\nats[\genericVecLen]}$ 
\item[\auditorLabel{auditor:sign}] $\forall \projIdx \in \nats[\numProjections], 
\genericSig[\projBucketsAugmented[\projIdx]] \assign 
\signingAlg[\pubKey](\projBucketsAugmented[\projIdx]) = 
\qrGen^{\left(\frac{1}{\blindingFactor_{\projIdx} \prod\limits_{\genericSetElmt \in 
\projBucketsAugmented[\projIdx]} 
		\genericSetElmt}\right)} \bmod \rsaMod$


\end{enumerate}

\smallskip\noindent\textbf{Client} ($\uploadFile(\clientData, \clientVec = 
\tupleDefn{\clientVec[1], 
\dots, 
	\clientVec[\genericVecLen]} \assign
\contentHash(\clientData), \setDefn{\genericSig[\projBucketsAugmented[\projIdx]]}_{\projIdx 
\in 
\nats[\numProjections]}, \pubKey)$)


\begin{enumerate}[leftmargin = 1.5em]
\item[\clientLabel{client:partition}] $\forall \projIdx \in \nats[\numProjections], 
\cInputSetAugmented[\projIdx] 
\assign \setDefn{\primeRO(\projIdx, \clientVec[\vecIdx], \vecIdx, 
\projFn[\projIdx](\clientVec))}_{\vecIdx 
	\in \nats[\genericVecLen]}$ 	 	
\item[\clientLabel{client:sign}] $\forall \projIdx \in \nats[\numProjections], \forall \vecIdx \in 
\nats[\genericVecLen], 
\genericSig[\projBucketsAugmented[\projIdx]] \sigMinus[\pubKey] \clientVec[\vecIdx] = 
\genericSig[\projBucketsAugmented[\projIdx]]^{\left(\clientVec[\vecIdx] \right)}  
\bmod \rsaMod$
\item[\clientLabel{client:randomize}] $\forall \projIdx \in \nats[\numProjections], \forall 
\vecIdx \in 
\nats[\genericVecLen], \sigPair[\projIdx, \vecIdx] \assign 
\randomize[\pubKey](\genericSig[\projBucketsAugmented[\projIdx]] \sigMinus[\pubKey] 
\clientVec[\vecIdx]) = 
\tupleDefn{\qrGen^{\left(\frac{\genericRand[\projIdx,\vecIdx] \cdot 
			\primeRO(\projIdx, \clientVec[\vecIdx], \vecIdx, 
			\projFn[\projIdx](\clientVec))}{\blindedBlocklistAccSetFuzzy{\projIdx}}\right)} 
	\bmod \rsaMod ;\ro(\qrGen^{\genericRand[\projIdx,\vecIdx]})}$
\item[\clientLabel{}] $\genericSymmKey \getsr \encryptKeySpace$
\item[\clientLabel{}] $\cOutputFile \assign \encrypt[\genericSymmKey](\clientData)$
\item[\clientLabel{client:share}] $\forall \projIdx \in \nats[\numProjections]$, 
$\secretShare{\projIdx, 1}, 
\dots, \secretShare{\projIdx, \genericVecLen} 
\assign 
\secretSharingScheme(\genericSymmKey)$
\item[\clientNoLabel] $\forall \projIdx \in \nats[\numProjections], \forall \vecIdx \in 
\nats[\genericVecLen], \cOutputSig[\projIdx][\vecIdx] \assign
\sigPair[\projIdx, \vecIdx][1]$
\item[\clientLabel{client:dhf}] $\forall \projIdx \in \nats[\numProjections], \forall \vecIdx \in 
\nats[\genericVecLen], \cOutputEnc[\projIdx][\vecIdx] \assign 
(\dhf[\dhfKey](\genericRandAlt[\projIdx, \vecIdx]), \secretShare{\projIdx, \vecIdx}) 
\oplus \sigPair[\projIdx,\vecIdx][2]$ where $\genericRandAlt[\projIdx, \vecIdx] \getsr 
\dhfDomain$
\item[\clientNoLabel] $\cOutputMeta \assign 
\setDefn{\cOutputSig[\projIdx][\vecIdx],\cOutputEnc[\projIdx][\vecIdx]}_{\forall \projIdx \in 
\nats[\numProjections], \forall \vecIdx \in 
	\nats[\genericVecLen]}$
	
\end{enumerate}

\smallskip\noindent\textbf{Server} ($\checkContent 
(\setDefn{\projBucketsAugmented[\projIdx]}_{\projIdx \in \numProjections}, 
\setDefn{\blindingFactor_{\projIdx}}_{\projIdx \in \numProjections}, \cOutputFile, 
\cOutputMeta, 
\setDefn{\genericSig[\projBucketsAugmented[\projIdx]]}_{\projIdx \in \numProjections}, 
\pubKey$)

\smallskip\noindent\algComment{Upon receiving $\tupleDefn{\cOutputFile, \cOutputMeta}$ from 
client, perform \stepsref{server:verify}{server:decrypt} for each $\projIdx \in 
\nats[\numProjections]$ and $\genericVec = \tupleDefn{\genericVec[1], \dots, 
\genericVec[\genericVecLen]}\in \blocklistEmbedded$}

\begin{enumerate}[leftmargin = 1.5em]

	\item[\serverLabel{server:verify}] 
$\forall \vecIdx \in \nats[\genericVecLen]$, $(\dhfOutput{\vecIdx}, \secretShare{\vecIdx}') 
\assign 
\ro(\cOutputSig[\projIdx][\vecIdx]\sigMinus[\pubKey]  
\projBucketsAugmented[\projIdx] 
	\setminus \setDefn{\genericVec[\vecIdx]}) \oplus \cOutputEnc[\projIdx][\vecIdx]$

\item[\serverLabel{server:dhf}]
$\matchedIndexSet \assign \dhf.\dhfDetect(\dhfOutput{1}, \dots, 
\dhfOutput{\genericVecLen})$. 

Abort if $\dhf.\dhfDetect(\dhfOutput{1}, \dots, \dhfOutput{\genericVecLen}) = \dhfFail$

\item[\serverLabel{server:combine}] $\genericSymmKey' = 
\secretSharingScheme.\SSCombine(\setDefn{\secretShare{\ssIdx}'}_{\ssIdx \in 
	\matchedIndexSet})$

\item[\serverLabel{server:decrypt}] $\clientData = 
\decrypt[\genericSymmKey'](\cOutputFile)$ if $\genericSymmKey' = \genericSymmKey$

\end{enumerate}
\end{mdframed}




\begin{theorem}
	
	The auditable content moderation protocol for detecting fuzzy matches against a blocklist 
	when using a randomizable set-homomorphic signature scheme \sigScheme satisfies the security 
	properties defined in \secref{sec:model}, that is 
	
	\iffalse
	\enumListStart
	
	\item It satisfies unforgeability of the blocklist token 
	
	\begin{footnotesize}
		\begin{gather*}
		 \cprob{\bigg}{\begin{array}{@{}r@{}}
					\genericSigAlt =  \createBlocklist\left(\setOfAuditorPrivkeys, \blocklistEmbedded' \right) \\
					\land ~\blocklistEmbedded' \not \subseteq \blocklistEmbedded
			\end{array}}{\begin{array}[m]{@{}l@{}}
					\tupleDefn{\genericSig[\projBuckets[\projIdx]]}{\projIdx \in \nats[\numProjections]}
					\assign 
					\createBlocklist\left(\setOfAuditorPrivkeys, 
					\blocklistEmbedded\right) \\
					\genericSigAlt \assign 
					\adversary\left(\tupleDefn{\genericSig[\projBuckets[\projIdx]]}{\projIdx \in 
					\nats[\numProjections]}, \blocklistEmbedded, 
					\begin{array}{@{}r@{}} \setOfAuditorPubkeys, \\ \setOfDishonestAuditorPrivkeys
					\end{array}\right)
			\end{array}} \\ \le \ufcmaAdv
		\end{gather*}
	\end{footnotesize}
	
	\item It hides non-matching content from the server 


	
	\begin{footnotesize}
		\begin{gather*}
			\cprob{\Bigg}{\begin{array}{@{}r@{}}
					\begin{array}{@{}r@{}} \adversary\bigg( \setOfAuditorPubkeys, \blocklist, \cOutputFile, 
					\cOutputMeta, \\
					\tupleDefn{\genericSig[\projBuckets[\projIdx]]}{\projIdx \in \nats[\numProjections]},  
					 \setOfDishonestAuditorPrivkeys
					\end{array}\bigg)  = \indBit
			\end{array}}{\begin{array}[m]{@{}l@{}}
				\tupleDefn{\genericSig[\projBuckets[\projIdx]]}{\projIdx \in \nats[\numProjections]}
				\assign 
				\createBlocklist\left(\setOfAuditorPrivkeys, 
				\blocklistEmbedded\right)
					\\
					(\clientData[0],\clientData[1]) \assign \adversary[1](	
					\tupleDefn{\genericSig[\projBuckets[\projIdx]]}{\projIdx \in \nats[\numProjections]}, 
					\blocklistEmbedded, 
					\setOfDishonestAuditorPrivkeys) \land \\	
					\genericCInputElmt_0 \assign \contentHash(\clientData[0]) ~\land 
					\genericCInputElmt_1 \assign \contentHash(\clientData[1]) ~\land  \\ 
					\forall\genericBlocklistElmt \in \blocklistEmbedded, 
					\hammingDist{\genericBlocklistElmt}{\genericCInputElmt_0} > \threshold ~\land\\
					\forall\genericBlocklistElmt \in \blocklistEmbedded, 
					\hammingDist{\genericBlocklistElmt}{\genericCInputElmt_1} > \threshold\\
					\indBit \assign \setDefn{0,1}\\
					(\cOutputFile, \cOutputMeta) \assign \uploadFile\left(\clientData, \genericCInputElmt, 
					\tupleDefn{\genericSig[\projBuckets[\projIdx]]}{\projIdx \in \nats[\numProjections]}\right)  
					\land 	~\genericCInputElmt \notin \blocklist
			\end{array}} \\ \le 1/2 + \Advantage{\indcpaLabel}{\encryptionSchemeDefn}(\timeBound, 
			\oracleQueries')  + \numProjections \genericVecLen \indsrAdv 
		\end{gather*}
	\end{footnotesize}
	
	\item It preserves confidentiality of the blocklist elements
	
	\begin{footnotesize}
		\begin{gather*}
			\cprob{\Bigg}{\begin{array}{@{}r@{}}
					\adversary[2]\left(\tupleDefn{\genericSig_{\indBit, \projIdx}}{\projIdx \in 
					\nats[\numProjections]},  
					\setOfAuditorPubkeys, \adversaryState{1}\right) = \indBit
			\end{array}}{\begin{array}[m]{@{}l@{}}
					(\blocklist_0, \blocklist_1, \adversaryState{1}) \assign 
					\adversary[1]\left(\setOfAuditorPubkeys\right) 
					\\
					\indBit \getsr \setDefn{0,1} \\
					\tupleDefn{\genericSig_{\indBit, \projIdx}}{\projIdx \in \nats[\numProjections]} \assign  
					\createBlocklist\left(\setOfAuditorPrivkeys, 
					\blocklist_\indBit\right) 
			\end{array}} \\ \le 1/2 + \frac{\numProjections \log \rsaMod}{\rsaMod}
		\end{gather*}
		
	\end{footnotesize}
	
	
	
	\enumListEnd  
	\fi
	
	
\end{theorem}



\begin{proofsketch}
	For (1), note that if there are no collisions in the random oracle $\primeRO$, then the sets 
	$\projBucketsAugmented[1], \dots, \projBucketsAugmented[\numProjections]$ do not have any 
	common items. As a result, the adversary does not gain any additional advantage in forging a 
	signature due to the blocklist partitioning. 

	For (2), note that each of the metadata items $\cOutputSig[\projIdx][\vecIdx]$ is generated 
	using an independent random coin $\genericRand[\projIdx, \vecIdx]$ and thus each of these 
	values are independent, uniformly distributed elements in $\qrGroup$. The values 
	$\cOutputEnc[\projIdx][\vecIdx]$ are encrypted with a one-pad where the pad is the output of a 
	random oracle on a random element in $\qrGroup$. The adversary breaks the one-time pad if: i) 
	it has queried the random oracle as a guess, or ii) it is able to compute the one-time pad
	from $\cOutputEnc[\projIdx][\vecIdx]$, which means it violates the IND-SR property of the 
	signature scheme. Thus, the adversary's advantage in obtaining the secret shares of the key for 
	any of 	$\numProjections \genericVecLen$ values is upper bounded by $\numProjections 
	\genericVecLen \indsrAdv$. Finally, since both the secret shares and the DHF outputs are 
	identically distributed to 
	uniformly random element in their respective ranges, with less than $\threshold$ correctly 
	decrypted shares (and DHF outputs), the adversary does not learn any information about the key. 
	Without the key, the adversary obtains the data from the ciphertext with advantage at most 
	$\Advantage{\indcpaLabel}{\encryptionSchemeDefn}(\timeBound, 
	\oracleQueries') $
	
	For (3), note that each of the signatures $\tupleDefn{\genericSig[\projIdx]}{\projIdx \in 
	\nats[\numProjections]}$ is blinded with an independent random prime number 
	$\blindingFactor_1, \dots, \blindingFactor_{\numProjections}$, and the sets have no common 
	elements. So by a union bound over the probability to break confidentiality of any of the blocklists, 
	we get the result. 
\end{proofsketch}


\myparagraph{Asymptotic costs}
%
The total size of the signatures is $2\numProjections$ \qrGroup elements. The client only needs to 
store these signatures. The communication volume between the client and the server 
includes the metadata which has $\numProjections \cdot \genericVecLen$ tuples, and each tuple 
contains one \qrGroup element, and a random oracle output with $\bigO(\secParam)$ 
bits. Also, the client performs $2\numProjections \genericVecLen$
exponentiations. The server's compute cost is $\bigO(\numProjections \genericVecLen \blocklistSize 
\log \genericVecLen \blocklistSize )$ exponentiations and $\bigO(\blocklistSize 
\numProjections)$ DHF detections. As discussed before, the server compute can be performed 
offline, and therefore does not impact the latency of file uploads. 










\iffalse

\setcounter{auditorLineNmbr}{0}
\setcounter{clientLineNmbr}{0}
\setcounter{serverLineNmbr}{0}


\begin{mdframed}[frametitle={Delegated Fuzzy PSI}, frametitlealignment=\centering, 
	font=\small]
	
	\smallskip\noindent\textbf{Parameters:} A randomizable set homomorphic signature scheme 
	$\sigScheme \assign \genericSigSchemeDefn$ with signature space $\sigOutputSpace$; an 
	IND-CPA secure encryption scheme $\encryptionSchemeDefn$ with key space 
	$\encryptKeySpace$; set of uniformly sampled functions $\projFn[1], \dots, 
	\projFn[\numProjections]: \setDefn{0,1}^{\genericVecLen} \rightarrow \projRange$ for a 
	locality-preserving partition over Hamming vectors; an embedding function 
	$\embed: \metricSpace \rightarrow \setDefn{0,1}^{\genericVecLen}$; 
	random oracles $\primeRO: 
	\setDefn{0,1} \times \nats[\genericVecLen] \times \projRange \rightarrow \primes$;  
	a $(\scalingFn(\threshold), \genericVecLen - \scalingFn(\threshold))$ detectable hash 
	function $\dhf: \dhfKeySpace \times \dhfDomain \rightarrow \dhfRange$;
	a secret-sharing scheme $\secretSharingScheme: \encryptKeySpace \rightarrow 
	\secretShareOutputSpace^{\genericVecLen}$; random oracle $\ro: 
	\sigOutputSpace \rightarrow \setDefn{0,1}^{\setSize{\dhfRange} + 
		\setSize{\secretShareOutputSpace}}$;
	
	\smallskip\noindent\textbf{Inputs:}
	\begin{itemize}[leftmargin=0.5em]
		\item The auditor \auditor has a blocklist \blocklist, the public key for the signature scheme 
		$\pubKey \assign \tupleDefn{\rsaMod, \qrGen}$ where $\rsaMod = \safePrime[1] \safePrime[2]$ 
		and $\qrGen$ is a random generator of the quadratic residue group \qrGroup of 
		\integersModNInv , and the corresponding private key $\privKey \assign \safePrime[1], 
		\safePrime[2]$  
		\item 	Client has set $\cInputSet = (\keySetLabel \times \valSetLabel) \subseteq \metricSpace 
		\times \setDefn{0,1}^{\ast}$ and the public key $\pubKey$ of the signature scheme
		\item Server has the set $\blocklist$ and the public key $\pubKey$ 	
	\end{itemize}
	
	
	\smallskip\noindent\textbf{Auditor}
	%
	\begin{enumerate}
		\item[\auditorLabel{auditor:embed}] $\blocklistEmbedded \assign \setDefn{
			\embed(\genericBlocklistElmt): \genericBlocklistElmt \in \blocklist }$ 
		\item[\auditorNoLabel] $\forall \projIdx \in \nats[\numProjections], 
		\projBuckets[\projIdx] \assign \emptyset$
		\item[\auditorLabel{auditor:partition}] $\forall \projIdx \in \nats[\numProjections], \forall 
		\genericVec \in \blocklistEmbedded, 
		\projBuckets[\projIdx] =  \projBuckets[\projIdx] \cup \projFn[\projIdx](\genericVec) ~\mbox{if}~ 
		\projFn[\projIdx](\genericVec) \notin \projBuckets[\projIdx]$ 
		\item[\auditorLabel{auditor:augment}] $\forall \projIdx \in \nats[\numProjections],
		\projBucketsAugmented[\projIdx] \assign \setDefn{\primeRO(\genericVec[\vecIdx], \vecIdx, 
			\projFn[\projIdx](\genericVec))}_{\genericVec \in \blocklistEmbedded, \vecIdx \in 
			\nats[\genericVecLen]}$ 
		\item[\auditorLabel{auditor:sign}] $\forall \projIdx, 
		\genericSig[\projBucketsAugmented[\projIdx]] \assign 
		\signingAlg[\pubKey](\projBucketsAugmented[\projIdx]) = 
		\qrGen^{\left(\frac{1}{\prod\limits_{\genericSetElmt \in \projBucketsAugmented[\projIdx]} 
				\genericSetElmt}\right)} \bmod \rsaMod$
	\end{enumerate}
	
	\smallskip\noindent\textbf{Client}
	%
	\begin{enumerate}
		\item[\clientLabel{client:embed}] $\cInputSetEmbedded \assign \setDefn{
			\embed(\genericCInputElmt.\keyLabel): \genericCInputElmt.\keyLabel \in \cInputKeySet}$
		\item[] \textit{/* Repeat \stepsref{client:partition}{client:dhf} for all $\clientVec \in 
			\cInputSetEmbedded$ */}
		\item[\clientLabel{client:partition}] $\forall \projIdx \in \nats[\numProjections], 
		\cInputSetAugmented[\projIdx] 
		\assign \setDefn{\primeRO(\clientVec[\vecIdx], \vecIdx, \projFn[\projIdx](\clientVec))}_{\vecIdx 
			\in \nats[\genericVecLen]}$ 	 	
		\item[\clientLabel{client:sign}] $\forall \projIdx \in \nats[\numProjections], \forall \vecIdx \in 
		\nats[\genericVecLen], 
		\genericSig[\projBucketsAugmented[\projIdx]\setminus\clientVec[\vecIdx]] = 
		\genericSig[\projBucketsAugmented[\projIdx]] \sigMinus[\pubKey] \clientVec[\vecIdx] = 
		\genericSig[\projBucketsAugmented[\projIdx]]^{\left(\clientVec[\vecIdx] \right)}  
		\bmod \rsaMod$
		\item[\clientLabel{client:randomize}] $\forall \projIdx \in \nats[\numProjections], \forall \vecIdx \in 
		\nats[\genericVecLen], \sigPair[\projIdx, \vecIdx] \assign 
		\randomize[\pubKey](\genericSig[\projBucketsAugmented[\projIdx]\setminus\clientVec[\vecIdx]]
		) = 
		\tupleDefn{\genericSig[\projBucketsAugmented[\projIdx]\setminus\clientVec[\vecIdx]]^{\genericRand[\projIdx,\vecIdx]},
			\ro(\qrGen^{\genericRand[\projIdx,\vecIdx]})}$
		\item[\clientLabel{}] $\genericSymmKey \getsr \encryptKeySpace$
		\item[\clientLabel{client:share}] $\forall \projIdx \in \nats[\numProjections]$, 
		$\secretShare{\projIdx, 1}, 
		\dots, \secretShare{\projIdx, \genericVecLen} 
		\assign 
		\secretSharingScheme(\genericSymmKey)$
		\item[\clientNoLabel] $\forall \projIdx \in \nats[\numProjections], \forall \vecIdx \in 
		\nats[\genericVecLen], \cOutputSig[\projIdx][\vecIdx] \assign
		\sigPair[\projIdx, \vecIdx][1]$
		\item[\clientLabel{client:dhf}] $\forall \projIdx \in \nats[\numProjections], \forall \vecIdx \in 
		\nats[\genericVecLen], \cOutputEnc[\projIdx][\vecIdx] \assign 
		(\dhf[\dhfKey](\genericRandAlt[\projIdx, \vecIdx]), \secretShare{\projIdx, \vecIdx}) 
		\oplus \sigPair[\projIdx,\vecIdx][2]$ where $\genericRandAlt[\projIdx, \vecIdx] \getsr 
		\dhfDomain$
	\end{enumerate}
	
	\smallskip\noindent\textbf{Server}
	
	\noindent\textit{/* Upon receiving $\setDefn{\cOutputSig[\projIdx][\vecIdx], 
			\cOutputEnc[\projIdx][\vecIdx]}_{\projIdx \in \nats[\numProjections], \vecIdx \in 
			\nats[\genericVecLen]}$ and $\cOutputVal \assign 
		\encrypt[\genericSymmKey](\genericCInputElmt.\valueLabel)$ for each of $\genericCInputElmt 
		\in \cInputSet$, repeat \stepsref{server:verify}{server:decrypt} for $\genericVec \in 
		\blocklistEmbedded, \projIdx \in \nats[\numProjections]$ */}
	
	
	
	\begin{enumerate}
		
		
		\item[\serverLabel{server:verify}] 
		$\forall \vecIdx \in \nats[\genericVecLen]$, $(\dhfOutput{\vecIdx}, \secretShare{\vecIdx}') 
		\assign 
		\ro(\cOutputSig[\projIdx][\vecIdx]\sigMinus[\pubKey]  
		\setDefn{\projBucketsAugmented[\projIdx] 
			\setminus \genericVec[\vecIdx]}) \oplus \cOutputEnc[\projIdx][\vecIdx]$
		
		\item[\serverLabel{server:dhf}]
		$\matchedIndexSet \assign \dhf.\dhfDetect(\dhfOutput{1}, \dots, \dhfOutput{\genericVecLen})$. 
		
		Abort if $\dhf.\dhfDetect(\dhfOutput{1}, \dots, \dhfOutput{\genericVecLen}) = \dhfFail$
		
		\item[\serverLabel{server:combine}] $\genericSymmKey' = 
		\secretSharingScheme.\SSCombine(\setDefn{\secretShare{\ssIdx}'}_{\ssIdx \in 
			\matchedIndexSet})$
		
		\item[\serverLabel{server:decrypt}] $\genericCInputElmt.\valueLabel = 
		\decrypt[\genericSymmKey'](\cOutputVal)$ if $\genericSymmKey' = \genericSymmKey$
		
	\end{enumerate}
	
\end{mdframed}
\fi